\honest{Reductive Reference}{Eliezer Yudkowsky}{2008}

The reductionist thesis, ``as I formulate it,'' is that minds use a map with many levels, atoms and quarks, heat and kinetic energy, while reality seems to have one. Carnot studied heat as a fluid long before anyone reduced it to motion, and that was efficient. Reality, by contrast, seems to compute everything at the bottom level. So in what sense is ``I am wearing socks'' true, when there are only quarks?

``The classic definition'' of truth, I say, is Tarski's: ``snow is white'' is true if and only if snow is white. Truth is ``an ideal comparison between a belief and reality,'' and you test it with telescopes, not introspection. \nb{Rival theories of truth accept Tarski's sentence too, so it does not support the comparison view over them.} ``This is the point missed by the postmodernist folks,'' I add, and I name none of them. I admit that reading an experiment depends on prior beliefs, and leave that for another post. Trusting the telescope because of what I believe about telescopes is, I say, not circular logic but ``reflective coherence.''

No one can compare a pattern of neurons with a snowflake quark by quark. ``It seems to me'' that a word like ``snow'' is a promissory note: a pointer to whatever simple boundary an ideal interpreter would draw around the stable cause of the white stuff I keep seeing. Putnam's Twin Earth, I say, ``helps to highlight this issue.'' \nb{The promissory note is close to the conclusion Putnam drew from Twin Earth, semantic externalism: the liquid Twin Earthers call ``water'' is not water, because the word refers to the stuff around the speaker, whatever its structure. The post does not say that its view is Putnam's.}

My beliefs about my beliefs are promissory notes too. ``So, rather than calling it circular, I prefer to call it self-consistent.'' \nb{Yudkowsky defended such loops at length a few months later, in ``Where Recursive Justification Hits Bottom''.} It is unnerving, but hard to see how it could be otherwise.

So reality has no fundamental hands, but I still have hands; if I lost them, reality would soon be in the state described as ``Eliezer screams as blood jets out of his wrist stumps.'' Higher levels are not false or empty: they exist implicitly in physics. \nb{Reduction without elimination is the usual view among philosophers of science: in the Stanford Encyclopedia's words, ``most reductive views are realist.''}
